\section{Online presence and threat intelligence}
\label{sec:online}

APT29 repeatedly uses services its targets already trust.

HAMMERTOSS, documented by FireEye in 2015, generated a new Twitter handle each day and read instructions from a tweet that pointed to an image on a site such as GitHub. The hashtag carried the offset of the hidden data and part of the key. Stolen data went to cloud storage accounts the operators created~\cite{fireeye-hammertoss}. A 2021 advisory describes infrastructure bought with false identities and cryptocurrency, often through temporary email accounts and resellers, including personas using email services on \texttt{cock[.]li} or related domains~\cite{aa21-116a}. No source reviewed here names a specific wallet or payment, so wallet attribution remains unverified.

The group also borrows legitimate mailing infrastructure. Microsoft reported that NOBELIUM used a compromised USAID Constant Contact account to send phishing mail to about 3,000 accounts at more than 150 organisations, with links through Constant Contact's redirector~\cite{ms-nobelium}. Polish authorities described embassy-themed emails that led to a compromised website and an ENVYSCOUT script using HTML smuggling~\cite{skw-certpolska}. Unit~42 documented a reused car-sale flyer sent to diplomatic missions in Kyiv, with payload control over the Microsoft Graph and Dropbox APIs~\cite{unit42}. In 2023, Microsoft reported Teams messages from new \texttt{onmicrosoft.com} domains created inside compromised tenants, which asked users to enter a code in the Microsoft Authenticator prompt~\cite{ms-teams}.

Every case relies on a service the target already allows, so blocking the actor's domains helps less than noticing when a trusted service delivers an untrusted request.
